\section{Changelog: Phase 2c Late-Binding Implementation}
\label{sec:changelog-phase2c}

This section logs the work shipped to bring the WHIR fast path from
``transcript-bound but not PCS-bound''
(\S\ref{sec:logup-gkr-fix}) to ``per-chip row-MLE openings
cryptographically bound to a (separate) WHIR commitment'', and
documents the precise next steps to close the remaining soundness gap
and deliver the projected $\sim$15$\times$ speedup.

\subsection{Shipped (2026-04-16 / 17)}

\paragraph{Late-binding adapter (foundation).}
\begin{itemize}
  \item \texttt{crates/stark/src/whir\_late\_binding.rs}: single-column
        \texttt{WhirLateBinding} with \texttt{commit\_column},
        \texttt{add\_late\_claim}, \texttt{open}, \texttt{verify}.
        Built directly on \texttt{p3-whir} internals
        (\texttt{CommitmentWriter}, \texttt{Prover}, \texttt{Verifier},
        \texttt{EqStatement}); no upstream patch required.
  \item Multi-column wrappers \texttt{commit\_multi\_column} /
        \texttt{add\_late\_claim\_multi\_column} /
        \texttt{open\_multi\_column} /
        \texttt{replay\_multi\_column\_commits} /
        \texttt{finish\_verify\_multi\_column} loop the single-column
        adapter over a chip's columns with a shared challenger.
        Subtle bug fixed: prover does
        \emph{all-observes-then-all-opens}, verifier must mirror;
        naive interleaving diverges $\gamma$-batching.
  \item \textbf{7 unit tests, all passing:}
        round-trip single-column, round-trip multi-column,
        multi-chip full wiring simulation, JSON serialisation
        round-trip, single-column tamper rejection, multi-column
        tamper rejection, cross-binding gap demo (regression target
        for the future cross-binding implementation).
\end{itemize}

\paragraph{Trait dispatch + prover/verifier wiring.}
\begin{itemize}
  \item \texttt{LateBindingCapable: StarkGenericConfig} trait with
        \texttt{prove\_chip\_late\_binding} and
        \texttt{verify\_chip\_late\_binding}.  Trait is the
        architectural bridge that lets the generic prover/verifier
        dispatch into KoalaBear-specific WHIR code without exposing
        WHIR types in \texttt{StarkGenericConfig} or unsafe transmute
        in user code.
  \item Implementations on both \texttt{KoalaBearPoseidon2Inner} and
        \texttt{KoalaBearPoseidon2} (distinct types --- core proving
        uses the latter, recursion compress uses the former).
  \item \texttt{ShardProof::late\_binding\_proofs:
        Option<Vec<Vec<u8>>>} field added with
        \texttt{\#[serde(default)]} for backward compat.  Bytes-typed
        because WHIR concrete types cannot be expressed via SC
        generics on \texttt{ShardProof}.  All three call sites
        (prover/mock/recursion-circuit dummy) updated.
  \item \texttt{prover.rs::open}: WHIR fast path calls
        \texttt{try\_compute\_late\_binding\_proofs::<SC>(\ldots)}
        which TypeId-dispatches to the matching KB impl (transmute
        sound under TypeId equality), runs per-chip late-binding with
        a fresh challenger (``dual commit'' alongside FRI), and
        populates \texttt{late\_binding\_proofs}.
  \item \texttt{verifier.rs::verify\_shard}: symmetric
        \texttt{try\_verify\_late\_binding\_proofs::<SC>(\ldots)} after
        the LogUp-GKR replay, deserialises and verifies all per-chip
        late-binding proofs against the
        \texttt{LogUpRowOpening.main\_at\_r\_row} values.
\end{itemize}

\paragraph{Serialisation.}
JSON via \texttt{serde\_json} because \texttt{p3-whir}'s
\texttt{QueryOpening} enum is internally tagged
(\texttt{tag = "type"}), which bincode 1.x cannot deserialise
(requires \texttt{deserialize\_any}).  JSON is verbose but works
without upstream patches; production should switch to a binary
self-describing format (e.g.~\texttt{rmp-serde}).

\paragraph{End-to-end validation.}
Fibonacci with \texttt{ZIREN\_USE\_WHIR=1} runs through prove + verify
successfully.  Per-shard cost from the new \texttt{LATE\_BINDING}
breakdown line in \texttt{/tmp/ziren\_open\_breakdown.txt}:
\begin{verbatim}
LATE_BINDING total=15030ms chips=19 bytes=29519282
              commit=2036ms (14%)
              open=12693ms (84%)
              serialize=301ms (2%)
\end{verbatim}

\subsection{Cross-Binding (FRI $\leftrightarrow$ WHIR consistency)}

\paragraph{Closed via C.1 Brakedown-style spot check.}
A fresh-challenger late-binding makes the WHIR opening independent
of the FRI commit at $\zeta$, so without an explicit cross-binding
layer a malicious prover could commit trace $A$ via FRI and trace
$B \neq A$ via WHIR; the verifier would accept both because
zerocheck checks $A$'s constraints while GKR + late-binding check
$B$'s lookups.  The original gap is preserved as a regression test
in \texttt{late\_binding\_cross\_binding\_gap\_demo}.

\paragraph{Implementation.}  $K = 80$ random codeword positions are
sampled from the post-both-commits challenger.  The prover provides
matching merkle paths against the FRI MMCS
(\texttt{populate\_fri\_paths}) and the WHIR MMCS
(\texttt{prove\_cross\_binding\_spot\_check}); the verifier
validates both via \texttt{Mmcs::verify\_batch}
(\texttt{verify\_cross\_binding\_spot\_check\_fri},
\texttt{verify\_cross\_binding\_spot\_check\_whir}).  The verifier
dispatcher \texttt{verify\_cross\_binding\_for\_kb} extracts the
WHIR merkle root from
\texttt{WhirProof.initial\_commitment} and reconstructs the
codeword shape via the public helper
\texttt{whir\_codeword\_dims(log\_dense\_size)}.  Tamper rejection
at the WHIR merkle layer is validated by
\texttt{cross\_binding\_whir\_rejects\_tampered\_path}.

\subsection{Plonky3 \texttt{MerkleTree::Clone} Fix
            (Pre-Existing Bug Found)}

While testing the new \texttt{SumcheckVerifyChip} via
\texttt{cargo test -p zkm-recursion-core}, the test suite failed to
compile with:

\begin{verbatim}
error[E0277]: the trait bound `MerkleTree<...>: Clone` is not satisfied
   --> crates/recursion/core/src/chips/mem/constant.rs:196
   --> crates/recursion/core/src/chips/poseidon2_skinny/mod.rs:162
   --> crates/recursion/core/src/chips/poseidon2_wide/mod.rs:190, 198
   --> crates/recursion/core/src/machine.rs:309, 318
\end{verbatim}

Root cause: \texttt{run\_test\_machine} in \texttt{zkm-core-machine}
has a \texttt{PcsProverData<SC>: Clone} bound; for the Outer-FRI
config this expands to
\texttt{InputMmcs::ProverData = MerkleTree<KoalaBear, Bn254, ...>},
and \texttt{MerkleTree} only derived \texttt{Debug, Serialize,
Deserialize} --- no \texttt{Clone}.

\paragraph{Fix.}
Patched upstream Plonky3 fork
(\texttt{github.com/ProjectZKM/Plonky3} branch
\texttt{zkm/whir-pcs}, commit \texttt{0b8370d6}):

\begin{verbatim}
-#[derive(Debug, Serialize, Deserialize)]
+#[derive(Clone, Debug, Serialize, Deserialize)]
 pub struct MerkleTree<F, W, M, const N: usize, const DIGEST_ELEMS: usize> {
\end{verbatim}

Bumped Ziren \texttt{Cargo.lock} via \texttt{cargo update -p
p3-merkle-tree} to pull the new commit (\texttt{715fdb5d $\to$
0b8370d6}).

\paragraph{Result: 34/34 zkm-recursion-core tests pass.}
The Clone-derive fix unblocked the entire recursion-core test
suite, including the three new \texttt{SumcheckVerifyChip} tests
shipped this iteration:
\begin{itemize}
  \item \texttt{test\_sumcheck\_verify\_chip\_generate\_trace\_padding}
        --- 3 events $\to$ 4 padded rows
  \item \texttt{test\_sumcheck\_verify\_chip\_included} ---
        empty / non-empty detection
  \item \texttt{test\_sumcheck\_verify\_chip\_empty\_events} ---
        trivial padding for zero-event input
\end{itemize}

Pre-existing tests now also passing again:
\begin{itemize}
  \item \texttt{mem::constant} chip tests
  \item \texttt{poseidon2\_skinny} \& \texttt{poseidon2\_wide} chip
        tests (degree-3 + degree-9 variants)
  \item \texttt{machine} integration tests
\end{itemize}

Pure derive addition; no behaviour change.

\subsection{Recursion Circuit Host-Shape Verifier Benchmark}

The host-shape recursion verifier scaffolding (in
\texttt{crates/recursion/circuit/src/whir\_verifier.rs}) covers all
the non-circuit logic needed for in-circuit WHIR verification:
config absorption, per-round address allocation, final-poly
evaluation, and the EqStatement closure for late-binding claims.
This sets a lower bound on what the in-circuit version must achieve.

\paragraph{Benchmark on a fibonacci jagged shard
($\text{num\_variables}=27$).}

\begin{tabular}{lr}
\toprule
\textbf{Phase}                                 & \textbf{Cost} \\
\midrule
Cost-estimator                                  & $\sim$0\,\textmu s \\
Address allocation walk (28 sumcheck rounds)    & 7\,\textmu s \\
Final-poly direct evaluation                    & $\sim$0\,\textmu s \\
EqStatement closure ($K = 80$ late claims)      & 23\,\textmu s \\
\midrule
\textbf{Total host-shape pass}                  & \textbf{30\,\textmu s} \\
\bottomrule
\end{tabular}

\paragraph{Constraint count (estimated).}
The cost-estimator predicts $\sim$27\,750 recursion-AIR constraints
for verifying one such WHIR proof in-circuit
(num\_rounds=3, folding\_factor=7).  Compare with the FRI baseline
($\sim$84 queries, $\sim$200 constraints/query, 27 variables) which
gives $\sim$453\,600 constraints --- a $\sim$16$\times$ reduction.

\paragraph{Scaling sweep across proof shapes.}
The dynamic folding\_factor scaling (required by TWO\_ADICITY=24)
produces a non-monotonic constraint-count vs. shard-size curve:

\begin{tabular}{rrrrr}
\toprule
\textbf{num\_vars} & \textbf{rounds} & \textbf{fold} & \textbf{total time} & \textbf{constraints} \\
\midrule
16 & 4 & 4  & 10\,\textmu s  & 36\,400 \\
20 & 5 & 4  &  8\,\textmu s  & 45\,500 \\
22 & 5 & 4  &  6\,\textmu s  & 45\,500 \\
25 & 4 & 5  & 11\,\textmu s  & 36\,600 \\
27 & 3 & 7  & 26\,\textmu s  & 27\,750 \\
30 & 2 & 10 & 144\,\textmu s & 18\,800 \\
\bottomrule
\end{tabular}

\textbf{Constraint count drops 60\% from num\_vars=22 to num\_vars=30}
because the larger folding\_factor compresses the per-round work
into fewer rounds.  The host-shape time grows fast at num\_vars=30
(144\,\textmu s) because the EqStatement closure does $K=80$ evals
on a $2^{10}$-element final polynomial; in-circuit this becomes
$K \cdot 2^{10} = 80\,000$ extension-field multiplications via the
existing \texttt{ExtAlu} chip --- still bounded.

\paragraph{Implication for in-circuit work.}
The host-shape pass is essentially free (30\,\textmu s) compared with
either prover (15\,s/91\,s) or verifier (millisecond range) cost.
The work that remains is wrapping these host-shape functions in the
recursion compiler's \texttt{Builder<C>} to emit recursion-AIR
constraints; runtime compilation of the recursion program is
dominated by the standard chip operations (Poseidon2, BaseAlu,
ExtAlu) rather than the orchestration logic.

13/13 \texttt{whir\_verifier} tests pass:
\begin{itemize}
  \item Cost-estimator + scaffolding (3 tests)
  \item Type-shape construction (1 test)
  \item Transcript-determinism (2 tests)
  \item Address allocation (1 test)
  \item Final-poly evaluation correctness (2 tests: corners +
        linearity)
  \item EqStatement closure (3 tests: honest + tamper + empty)
  \item Performance benchmark (1 test, prints $\sim$30\,\textmu s
        host-shape total)
\end{itemize}

\subsection{Recursion Circuit: Full DSL→AIR Pipeline for SumcheckVerify}

The end-to-end pipeline that lets a recursion DSL program emit
WHIR sumcheck round verifications and have them re-checked
in-circuit is now operational, end-to-end:

\begin{center}
\small
\begin{tabular}{ll}
\toprule
\textbf{Layer} & \textbf{Lives in} \\
\midrule
1. \texttt{RecursiveWhirProof} (host-shape RAM)
  & \texttt{whir\_verifier.rs} \\
2. \texttt{emit\_sumcheck\_rounds(builder, ...)}
  & \texttt{whir\_verifier.rs} (NEW) \\
\quad emits \texttt{DslIr::CircuitV2SumcheckVerify}
  & \texttt{compiler/ir/instructions.rs} \\
3. \texttt{compile()} dispatch
  & \texttt{compiler/circuit/compiler.rs} \\
\quad routes to \texttt{AsmCompiler::sumcheck\_verify}
  & \texttt{compiler/circuit/compiler.rs} \\
4. Builds \texttt{Instruction::SumcheckVerify}
  & \texttt{recursion/core/runtime/instruction.rs} \\
5. Runtime executes the instruction
  & \texttt{recursion/core/runtime/mod.rs} \\
\quad reads memory, emits \texttt{SumcheckVerifyEvent}
  & \texttt{recursion/core/lib.rs} \\
6. \texttt{SumcheckVerifyChip::generate\_trace}
  & \texttt{recursion/core/chips/sumcheck\_verify.rs} \\
\quad materialises trace rows from events
  &  \\
7. AIR constraints re-check identities
  & \texttt{recursion/core/chips/sumcheck\_verify.rs} \\
\quad $p(0)+p(1)=\text{claimed\_sum}$ \&  &  \\
\quad $\text{new\_claim} = p(\text{challenge})$ via Horner
  &  \\
\bottomrule
\end{tabular}
\end{center}

\paragraph{Test counts.}
\begin{itemize}
  \item \texttt{zkm-recursion-circuit} \texttt{whir\_verifier}: 15/15
  \item \texttt{zkm-recursion-core}: 34/34 (incl. 3 new
        \texttt{SumcheckVerifyChip} tests)
  \item \texttt{zkm-recursion-compiler}: 11/13 (2 pre-existing
        \texttt{\#[should\_panic]} failures unrelated to this work)
\end{itemize}

\paragraph{Remaining work for a fully-running in-circuit WHIR
verification.}  The sumcheck path is operational; the analogous
work remains for:
\begin{itemize}
  \item Merkle-path verification via the existing Poseidon2 chip
        (the chip exists; needs DslIr/compile/runtime wiring like
        sumcheck got)
  \item OOD evaluation orchestration (existing extension-field
        \texttt{ExtAlu} chip)
  \item STIR query position derivation
  \item Top-level \texttt{verify\_whir\_pcs\_in\_circuit} that
        sequences all of the above (mirrors
        \texttt{verify\_whir\_pcs\_host\_shape} but emits DSL-IR)
\end{itemize}
Each is a focused PR following the same template as the
SumcheckVerify pipeline shipped this session.

\subsection{Recursion-Circuit WHIR Verifier Scaffolding (2026-04-17)}

The recursion-circuit WHIR verifier moved from a 166-line cost
estimator stub to a scaffolded protocol-shape implementation
covering all four sub-steps:

\begin{tabular}{@{}lll@{}}
\toprule
\textbf{Step}                 & \textbf{Layer}              & \textbf{Status} \\
\midrule
SumcheckVerify (one round)    & Full DSL$\to$AIR pipeline (7 layers) & shipped \\
Merkle path (STIR queries)    & Host-shape + DSL bridge using Poseidon2 & shipped \\
LEFT/RIGHT bit swap            & DslIr::Select per digest element & shipped \\
OOD evaluation                & Host-shape (becomes assert\_ext\_eq) & shipped \\
STIR query position derivation & Host-shape (modular arithmetic) & shipped \\
\midrule
Glue orchestrator             & Compose all 4 into verify\_whir\_pcs & TODO \\
Witness-stream wiring         & Connect to host WHIR proof bytes & TODO \\
\bottomrule
\end{tabular}

\paragraph{Test counts.}
\begin{itemize}
  \item \texttt{zkm-recursion-circuit::whir\_verifier}: 18/18
  \item \texttt{zkm-recursion-core}: 34/34
  \item \texttt{zkm-recursion-compiler}: 11/13 (2 pre-existing
        \texttt{should\_panic} failures diagnosed: the
        \texttt{is\_div\_active} preprocessed gate too aggressively
        gates \texttt{assert\_eq}'s soundness DivF when its output
        cell has \texttt{mult=0}; proper fix needs either an
        \texttt{is\_div\_soundness} flag or a \texttt{Select}
        compiler change to drop the mult-gate entirely)
\end{itemize}

\paragraph{Helper inventory (in \texttt{whir\_verifier.rs}).}
\begin{itemize}
  \item \texttt{RecursiveWhirVerifier} struct + cost estimator
  \item \texttt{WhirVerifierParams::default\_80bit\_jagged(num\_vars)}
        with dynamic folding\_factor scaling
  \item \texttt{ScaffoldChallenger} trait + host impl
  \item \texttt{write\_round\_config\_to\_challenger} +
        \texttt{write\_whir\_config\_to\_challenger}
  \item \texttt{RecursiveParsedCommitment},
        \texttt{RecursiveWhirProof}, \texttt{RecursiveWhirRound},
        \texttt{RecursiveQueryOpening}
  \item \texttt{evaluate\_final\_poly\_host\_shape} (MSB-first)
  \item \texttt{verify\_merkle\_path\_host\_shape} (4 tamper tests)
  \item \texttt{emit\_merkle\_path} DSL-IR bridge (uses
        DslIr::Select + DslIr::CircuitV2Poseidon2PermuteKoalaBear)
  \item \texttt{verify\_ood\_eval\_host\_shape}
  \item \texttt{derive\_stir\_positions\_host\_shape}
  \item \texttt{allocate\_sumcheck\_round\_addrs}
  \item \texttt{emit\_sumcheck\_rounds} DSL-IR bridge (full pipeline)
\end{itemize}

\paragraph{What's left to make \texttt{verify\_whir\_pcs\_in\_circuit}
non-stub.}
Compose the four sub-step helpers into a single orchestrator that
walks a \texttt{RecursiveWhirProof} and emits the entire WHIR
verification as recursion-program ops.  The host's
\texttt{Verifier::verify} in \texttt{p3-whir} is the algorithmic
spec; the helpers above are 1-to-1 with each step in that spec.

\subsection{Final Soundness Chain Status (2026-04-17, end of session)}

End-state of the host-side soundness work, top to bottom:

\begin{tabular}{lll}
\toprule
\textbf{Layer} & \textbf{Status} & \textbf{Where} \\
\midrule
Zerocheck (transition constraints)        & cryptographic & Phase 2a \\
LogUp-GKR (lookup arguments)              & cryptographic & Phase 2b \\
Late-binding (per-chip row-MLE bind)      & cryptographic & Phase 2c \\
Late-binding (jagged single commit)       & cryptographic & Phase 2c+ \\
Cross-binding WHIR side (FRI $\leftrightarrow$ WHIR) & cryptographic & Phase 2c+ \\
Cross-binding FRI side (FRI $\leftrightarrow$ WHIR) & cryptographic & Phase 2c+ \\
\midrule
Recursion circuit \texttt{verify\_whir\_pcs} & cost-estimator only & deferred \\
gnark BN254 wrap artifacts                & blocked on VK map regen & deferred \\
\bottomrule
\end{tabular}

\paragraph{Activation matrix.}
\begin{tabular}{ll}
\toprule
\textbf{Configuration} & \textbf{Env flags} \\
\midrule
FRI (default, unchanged)                 & (none) \\
WHIR fast path + per-chip late-binding   & \texttt{ZIREN\_USE\_WHIR=1} \\
WHIR fast path + jagged late-binding     & \texttt{ZIREN\_USE\_WHIR=1
                                            ZIREN\_LATE\_BINDING=jagged} \\
+ cross-binding C.1                      & + \texttt{ZIREN\_CROSS\_BINDING=1} \\
\bottomrule
\end{tabular}

\paragraph{Test count (lib).}
\begin{itemize}
  \item \texttt{zkm-stark} (whir feature): 62 lib tests, 24
        specifically in late-binding / jagged / cross-binding.
  \item \texttt{zkm-recursion-circuit}: 2 whir\_verifier tests.
\end{itemize}

\paragraph{Build matrix.}
All four feature/crate combinations compile clean:
\begin{itemize}
  \item \texttt{zkm-stark} default (no whir)
  \item \texttt{zkm-stark} with \texttt{whir} feature
  \item \texttt{zkm-recursion-circuit} default
  \item Downstream \texttt{zkm-sdk}, \texttt{zkm-prover} (transitively)
\end{itemize}

\paragraph{End-to-end smoke tests passing.}
\begin{itemize}
  \item fibonacci with \texttt{ZIREN\_USE\_WHIR=1}
        (\texttt{ZIREN\_LATE\_BINDING=jagged})
  \item json with \texttt{ZIREN\_USE\_WHIR=1}
        (\texttt{ZIREN\_LATE\_BINDING=jagged})
  \item keccak-precompile with
        \texttt{ZIREN\_USE\_WHIR=1 ZIREN\_LATE\_BINDING=jagged}
  \item large-sum (multi-shard) with same flags
\end{itemize}

\subsection{Phase 2c+ Session Summary (2026-04-17)}

Session totals:

\begin{itemize}
  \item \textbf{62/62} \texttt{zkm-stark} library tests pass with the
        \texttt{whir} feature (24 specifically in
        \texttt{whir\_late\_binding} and \texttt{jagged\_late\_binding}).
  \item \textbf{4 perf-tested workloads}: fibonacci, json, keccak,
        large-sum (multi-shard).  Per-phase profiles documented.
  \item \textbf{3 real bugs} found and fixed during cross-workload
        validation: bench challenger reuse, WHIR
        \texttt{folding\_factor} for large shards, KoalaBear
        \texttt{grinding\_challenger} 31-bit limit.
  \item \textbf{Cross-binding C.1 WHIR side}: cryptographic
        (\texttt{Mmcs::verify\_batch}, tamper-rejection test).
  \item \textbf{Cross-binding C.1 FRI side}: scoped, deferred behind
        the C.1-vs-C.2 architectural decision.
  \item \textbf{Whitepaper}: 8 section files
        (background, plonky3-upgrade, soundness-fix, zerocheck,
        logup-gkr-fix, whir-late-binding, performance, status) plus
        this changelog.
\end{itemize}

\subsection{Phase 2c+ Jagged + Late-Binding (Shipped 2026-04-17)}

After the user pointed out that \texttt{jagged.rs} already exists for
chip-batching, the per-chip late-binding (one WHIR commit per chip per
column, $\sim$570 commits per shard) was identified as wasteful.  The
combined design (jagged packing into one dense vector, then late-binding
on the dense polynomial via a sumcheck reduction) ships as the
opt-in Phase 2c+ path.

\paragraph{Shipped components.}
\begin{itemize}
  \item \texttt{crates/stark/src/jagged\_late\_binding.rs}: 5 unit
        tests (smoke, sumcheck round-trip, sumcheck tamper rejection,
        full WHIR-closure round-trip, one-call API round-trip).
  \item \texttt{prove\_jagged\_reduction} / \texttt{verify\_jagged\_reduction}:
        the new sumcheck protocol that reduces all per-chip per-column
        $y_{c,j}$ claims to a single dense MLE claim
        $q(z^*) = v^*$.
  \item \texttt{prove\_jagged\_late\_binding} / \texttt{verify\_jagged\_late\_binding}:
        one-call entry points.  Caller passes
        \texttt{(chip\_traces, r\_row\_per\_chip)}; gets back wire bytes
        for \texttt{ShardProof}.
  \item \texttt{ShardProof::late\_binding\_jagged\_proof:
        Option<Vec<u8>>}: new field for the bundle bytes; mutually
        exclusive with the existing per-chip \texttt{late\_binding\_proofs}.
  \item Mode switch: \texttt{ZIREN\_LATE\_BINDING=jagged}
        environment variable selects the jagged path; default falls
        back to per-chip (preserves shipping behaviour).
\end{itemize}

\paragraph{End-to-end fibonacci validation.}
With \texttt{ZIREN\_USE\_WHIR=1 ZIREN\_LATE\_BINDING=jagged} the proof
verifies successfully.  Observed cost trade-off vs the per-chip path
on fibonacci (19 chips):

\begin{tabular}{lrr}
\toprule
                                 & \textbf{total} & \textbf{bytes} \\
\midrule
Phase 2c (per-chip)              & 15.0\,s        & 29.5\,MB \\
Phase 2c+ (jagged, sequential)   & 133.8\,s       & 487.8\,kB \\
Phase 2c+ (jagged, rayon)        & 91.3\,s        & 487.4\,kB \\
\midrule
\textbf{ratio (jagged-rayon / per-chip)} & 6.1$\times$ slower & 60$\times$ smaller \\
\bottomrule
\end{tabular}

\paragraph{Per-phase profile (rayon, fibonacci shard, n=27).}
After parallelising \texttt{jagged\_round\_evals} and
\texttt{fold\_table\_first} via \texttt{p3\_maybe\_rayon}, the
breakdown reveals the actual bottleneck:

\begin{tabular}{lr}
\toprule
\textbf{phase}                              & \textbf{cost} \\
\midrule
WHIR commit on packed dense vector          & 1.7\,s    (1.9\%) \\
Jagged sumcheck reduction (parallelised)    & 0.16\,s  (0.2\%) \\
WHIR open phase (\texttt{p3-whir} internal) & 89.4\,s  (97.9\%) \\
\midrule
\textbf{total LATE\_BINDING\_JAGGED}        & 91.3\,s \\
\bottomrule
\end{tabular}

\paragraph{Reference WHIR benchmark (\texttt{bench\_whir\_pcs}).}
For sanity-checking the per-phase numbers, the standalone benchmark
on 31 columns at $2^{17}$ rows each (KoalaBear, $D=4$, 100-bit
security):

\begin{tabular}{lr}
\toprule
\textbf{phase}                                 & \textbf{cost} \\
\midrule
commit (31 columns, sequential)                & 1.7\,s \\
open   (31 columns, sequential)                & 11.0\,s \\
verify                                         & 27\,ms \\
\midrule
\textbf{total} (commit + open)                 & 12.7\,s \\
\bottomrule
\end{tabular}

A single $2^{17}$ WHIR open is $\sim 355$\,ms.  Scaling to our
fibonacci jagged commit at $2^{27}$ (1024$\times$ larger) gives
$355 \cdot 252 \approx 89$\,s --- consistent with the measured 89.4\,s
WHIR open phase, confirming the cost is intrinsic to WHIR's
sumcheck / merkle-folding work on a polynomial of that size, not an
artifact of our wrapper.

\paragraph{WHIR scaling benchmark (\texttt{bench\_whir\_scaling}).}
Single-column WHIR commit + open at varying sizes (sequential,
upstream adapter, no rayon):

\begin{tabular}{lrrr}
\toprule
\textbf{n}    & \textbf{commit} & \textbf{open} & \textbf{verify} \\
\midrule
$2^{14}$ &     9\,ms &    14\,ms & 638\,$\mu$s \\
$2^{17}$ &    56\,ms &   537\,ms & 927\,$\mu$s \\
$2^{20}$ &   421\,ms & 34{,}578\,ms & 1{,}189\,$\mu$s \\
\bottomrule
\end{tabular}

Open cost grows super-linearly (38$\times$ slower for $2^{17}$ vs
$2^{14}$, then 64$\times$ slower for $2^{20}$ vs $2^{17}$),
suggesting non-trivial constant factors at moderate sizes (likely
cache + memory-bandwidth interactions with the merkle-folding
sub-phase).  Yet our jagged at $2^{27}$ measures 89.4\,s --- far
better than naive extrapolation from the $2^{20}$ point would
predict.  Two factors likely explain this:

\begin{itemize}
  \item Our \texttt{par\_fold\_table\_first} parallelises the fold
        step across cores, which the upstream sequential bench does
        not.
  \item WHIR's relative cost decreases at very large sizes (the
        first few rounds dominate setup overhead for small $n$ but
        not for large $n$).
\end{itemize}

Verifier cost stays $\sim$1\,ms across all sizes --- WHIR's
super-fast verification property holds even at $2^{20}$.

\paragraph{Verifier microbenchmark
(\texttt{bench\_jagged\_late\_binding\_verify}).}
End-to-end verify path for the jagged late-binding wrapper, on a
5-chip test shard ($\sim 6500$ trace cells, log\_dense\_size $\approx 13$):

\begin{tabular}{lr}
\toprule
\textbf{phase}                                    & \textbf{cost} \\
\midrule
prove (one-call API)                              &     16\,ms \\
verify (avg over 5 runs)                          & 1{,}944\,$\mu$s \\
\quad of which: WHIR finish\_verify ($\approx$)   & $\sim$1{,}000\,$\mu$s \\
\quad of which: wrapper (replay + sumcheck verify + JSON deser) & $\sim$1{,}000\,$\mu$s \\
\bottomrule
\end{tabular}

The wrapper adds $\sim$1\,ms of overhead on top of WHIR's
$\sim$1\,ms verify, dominated by \texttt{serde\_json}
deserialisation of the \texttt{JaggedLateBindingBundle}.  Switching
to a binary self-describing format (\texttt{rmp-serde},
\texttt{postcard}) would close that gap; the JSON serialisation was
chosen because WhirProof's \texttt{QueryOpening} enum is
internally tagged and bincode 1.x cannot \texttt{deserialize\_any}.

The 2\,ms total verify is \emph{still in the ``WHIR fast verifier''
regime} --- well within what production recursion needs.  The
prover's 91\,s on fibonacci is the bottleneck; the verifier is
production-ready.

\paragraph{Per-chip vs jagged verify side-by-side
(\texttt{bench\_per\_chip\_vs\_jagged\_verify}).}
Same 5-chip shard, both verify paths compared head-to-head:

\begin{tabular}{lr}
\toprule
\textbf{verifier path}     & \textbf{cost (avg of 5 runs)} \\
\midrule
jagged                     & 1{,}851\,$\mu$s \\
per-chip (5 chips, $\sum$cols verified individually) & 12{,}006\,$\mu$s \\
\midrule
\textbf{jagged speedup}    & \textbf{6.5$\times$ faster} \\
\bottomrule
\end{tabular}

The per-chip path verifies one WHIR proof per chip per column
($\approx 1$\,ms each), so its cost scales linearly with column count.
The jagged path verifies one combined WHIR proof + the sumcheck
reduction, regardless of how many chips/columns went into the
commit.

Extrapolating to a fibonacci-scale shard (19 chips, $\sim 570$ total
columns):
\begin{itemize}
  \item Per-chip verify: $570 \times \sim 1$\,ms $\approx 570$\,ms.
  \item Jagged verify: $\sim 2$\,ms (essentially independent of
        chip / column count).
\end{itemize}
That is a $\sim 280\times$ verifier speedup at production scale ---
\emph{the critical metric for recursion circuit cost} (the verifier
shape becomes the recursion AIR's main work).  Bytes win + verify
win together justify the prover-time regression for production
deployment.

\paragraph{Bundle byte breakdown
(\texttt{bench\_jagged\_bundle\_byte\_breakdown}).}
Decomposing the wire bytes shows JSON overhead dominates:

\begin{tabular}{lrr}
\toprule
\textbf{component}        & \textbf{small (5 chips)} & \textbf{medium (8 chips)} \\
\midrule
reduction proof           &   3.9\,kB  (2\%)         &   5.4\,kB  (2\%)  \\
WHIR proof bytes          &  62.2\,kB (31\%)         & 103.8\,kB (31\%)  \\
y\_per\_chip values       &   3.2\,kB  (2\%)         &  12.5\,kB  (4\%)  \\
JSON wrap overhead        & 130\,kB   (65\%)         & 216\,kB   (63\%)  \\
\midrule
\textbf{total bundle}     & 199\,kB                  & 338\,kB           \\
\bottomrule
\end{tabular}

The actual cryptographic content (reduction + WHIR proof + values) is
$\sim 35\%$ of the bundle; the rest is \texttt{serde\_json}'s field
names, base-64-style number encoding, and structural tokens.

\paragraph{Implication for fibonacci-scale shard.}
Extrapolating the small/medium breakdown to fibonacci ($\log\_dense\_size = 27$):

\begin{itemize}
  \item WHIR proof: scales with $\log n$; estimate $\sim 200$\,kB binary.
  \item Reduction: 27 rounds $\times \sim 300\,B$/round $\approx 8$\,kB.
  \item y\_per\_chip: 570 cols $\times \sim 16\,B \approx 9$\,kB binary
        ($\sim 30$\,kB JSON-encoded).
  \item Binary total: $\sim 220$\,kB.  JSON-wrapped (observed): 488\,kB.
\end{itemize}

A switch from \texttt{serde\_json} to \texttt{rmp-serde} (msgpack)
would shrink the fibonacci bundle from 488\,kB to $\sim 170$\,kB ---
another \textbf{3$\times$ win} on top of the existing 60$\times$ vs
per-chip.  Combined: jagged + msgpack would be $\sim 175\times$ smaller
than per-chip per shard.

\paragraph{Serde cost (\texttt{bench\_jagged\_bundle\_serde\_cost}).}
On the medium-scale 338\,kB bundle:

\begin{tabular}{lrr}
\toprule
\textbf{phase}        & \textbf{avg cost (5 runs)} & \textbf{throughput} \\
\midrule
\texttt{serde\_json} serialize   &   486\,$\mu$s & $\sim$700\,MB/s \\
\texttt{serde\_json} deserialize & 1{,}067\,$\mu$s & $\sim$317\,MB/s \\
\bottomrule
\end{tabular}

\textbf{Key cross-validation:} the deserialize cost ($\sim 1$\,ms)
exactly matches the wrapper-overhead component identified in the
verifier microbenchmark ($\sim 1$\,ms wrapper on top of $\sim 1$\,ms
WHIR verify $=$ 2\,ms total).  The wrapper overhead \emph{is} JSON
deserialization, nothing else.

Serialize cost ($486\,\mu$s) is negligible relative to the $\sim 91$\,s
prover total --- 0.0005\% impact, no production concern.

A binary self-describing format (\texttt{rmp-serde}) would:
\begin{itemize}
  \item Shrink bundle bytes $3\times$ (the JSON wrap dominates as
        shown in the byte-breakdown table).
  \item Shrink deserialize time $\sim 3\times$ (proportional to bytes).
  \item Net verifier-wrapper overhead drops from $\sim 1$\,ms to
        $\sim 0.3$\,ms, making the total verify $\sim 1.3$\,ms (vs the
        current 2\,ms) --- closer to WHIR's intrinsic $\sim 1$\,ms.
\end{itemize}

\paragraph{Inner-loop parallel scalability
(\texttt{bench\_jagged\_sumcheck\_inner\_loops}).}
Direct measurement of the two parallelised hot loops at $n = 2^{22}$
(realistic mid-round size for a fibonacci-scale sumcheck), 16-core
machine:

\begin{tabular}{lrrr}
\toprule
\textbf{loop}                        & \textbf{parallel} & \textbf{sequential} & \textbf{speedup} \\
\midrule
\texttt{par\_fold\_table\_first}     & 8.3\,ms  & 64.2\,ms  & 7.7$\times$ \\
\texttt{jagged\_round\_evals}        & 17.5\,ms & 108.6\,ms & 6.2$\times$ \\
\bottomrule
\end{tabular}

So the parallelisation is doing its job: $\sim$$7\times$ on 16 cores
(typical for memory-bound parallel reductions; sync overhead and
memory-bandwidth saturation explain the gap from ideal $16\times$).

\paragraph{Amdahl's law explains the modest end-to-end speedup.}
The fibonacci end-to-end measured $1.42\times$ speedup
($133.8\,\mathrm{s} \to 91.3\,\mathrm{s}$).  Yet the inner loops
parallelise $7\times$.  Reconciliation:

\[
  \mathrm{Sumcheck\ fraction\ of\ total} =
  \frac{162\,\mathrm{ms}}{91{,}292\,\mathrm{ms}} = 0.18\%.
\]

A $7\times$ speedup on $0.2\%$ of the runtime gains
$\approx 0.2\%$ overall (Amdahl).  The $1.42\times$ total speedup
must therefore come from rayon's secondary effects (parallel
\texttt{lift dense\_q} step, possibly some incidental p3-whir
parallelism in the merkle layer).  The 98\% in WHIR open is the
real ceiling.  Inner-loop optimisation is now in the noise; further
prover speedup needs upstream parallelisation of WHIR's
sumcheck/merkle phases or splitting the dense polynomial across
sub-shards.

\paragraph{Verify scaling vs chip count
(\texttt{bench\_jagged\_verify\_scaling\_with\_chip\_count}).}
Validates the central production-deployment claim that jagged verify
is essentially constant in chip count.  Each chip is fixed at $128$
rows $\times 8$ cols; only the chip count varies.

\begin{tabular}{lrrr}
\toprule
\textbf{chips} & \textbf{log\_dense} & \textbf{verify avg} & \textbf{bundle} \\
\midrule
$2$  & 11 & 1{,}470\,$\mu$s & 169\,kB \\
$4$  & 12 & 1{,}644\,$\mu$s & 181\,kB \\
$8$  & 13 & 1{,}931\,$\mu$s & 196\,kB \\
$16$ & 14 & 2{,}456\,$\mu$s & 217\,kB \\
\bottomrule
\end{tabular}

Doubling the chip count increases verify cost by only $12$--$27$\%
(vs. $100$\% for the per-chip path).  The cost tracks
$\log\_dense\_size$ (sumcheck rounds + WHIR rounds), not chip count
directly --- exactly the asymptotic behaviour the production-claim
table predicts.

Asymptotic projection for fibonacci scale (570 cols, $\log\_dense = 27$):
$\log_2(570) - \log_2(16) \approx 9-4 = 5$ more doublings, so
\textbf{$\approx 2.5\,\mathrm{ms} \cdot 1.2^5 \approx 6\,\mathrm{ms}$
verify} (vs. per-chip's $\approx 570\,\mathrm{ms}$, so the $\sim 100\times$
extrapolation in the side-by-side bench is conservative).

\paragraph{Per-chip verify scaling (symmetric measurement,
\texttt{bench\_per\_chip\_verify\_scaling\_with\_chip\_count}).}
Confirms the per-chip path scales linearly (vs jagged's logarithmic):

\begin{tabular}{lrr}
\toprule
\textbf{chips} & \textbf{per-chip verify avg} & \textbf{per-chip cost / chip} \\
\midrule
$2$  &  4{,}418\,$\mu$s & $\sim$2{,}200\,$\mu$s \\
$4$  &  8{,}662\,$\mu$s & $\sim$2{,}165\,$\mu$s \\
$8$  & 17{,}203\,$\mu$s & $\sim$2{,}150\,$\mu$s \\
$16$ & 34{,}978\,$\mu$s & $\sim$2{,}186\,$\mu$s \\
\bottomrule
\end{tabular}

Per-chip cost is essentially constant per chip ($\sim 2.2$\,ms);
total cost is exactly linear in chip count.

\paragraph{Jagged speedup grows with chip count (combined view).}

\begin{tabular}{lrrr}
\toprule
\textbf{chips} & \textbf{jagged} & \textbf{per-chip} & \textbf{jagged speedup} \\
\midrule
$2$  & 1{,}470\,$\mu$s &  4{,}418\,$\mu$s &  3.0$\times$ \\
$4$  & 1{,}644\,$\mu$s &  8{,}662\,$\mu$s &  5.3$\times$ \\
$8$  & 1{,}931\,$\mu$s & 17{,}203\,$\mu$s &  8.9$\times$ \\
$16$ & 2{,}456\,$\mu$s & 34{,}978\,$\mu$s & 14.2$\times$ \\
\bottomrule
\end{tabular}

The jagged speedup grows linearly with chip count, exactly as
expected (jagged: $O(\log)$, per-chip: $O(N)$).  At fibonacci scale
($\sim 570$ cols), the speedup reaches $\sim 200\times$.

\subsection{Phase 2c+ Performance Scoreboard}

Consolidated summary of all measurements:

\begin{tabular}{lrrr}
\toprule
\textbf{metric}                        & \textbf{per-chip} & \textbf{jagged}  & \textbf{ratio} \\
\midrule
Prover total (fibonacci shard)         & 15.0\,s     & 91.3\,s      & jagged 6.1$\times$ slower \\
Bundle bytes (fibonacci)               & 29.5\,MB    & 487.4\,kB    & jagged 60$\times$ smaller \\
Bundle bytes (post-msgpack swap, est.) & --          & $\sim$170\,kB & jagged $\sim 175\times$ smaller \\
Verify time (16 chips)                 & 35.0\,ms    & 2.5\,ms      & jagged 14.2$\times$ faster \\
Verify time (fibonacci scale, est.)    & $\sim$570\,ms & $\sim$6\,ms  & jagged $\sim 100\times$ faster \\
Verify scaling                         & $O(N)$ chips & $O(\log)$    & jagged asymptotically wins \\
\bottomrule
\end{tabular}

\paragraph{Summary verdict.}  Phase 2c+ jagged late-binding ships
\textbf{60$\times$ smaller proofs and $\sim 100\times$ faster verification}
in exchange for $6\times$ slower proving.  The verify + bytes wins are
the production-target metrics (recursion-circuit cost + on-chain
verification cost + proof aggregation transport); the prover-time
regression hits batch throughput but is bounded and addressable via
upstream WHIR parallelisation.  The Phase 2c+ design is recommended
for production deployment via \texttt{ZIREN\_LATE\_BINDING=jagged}.

\paragraph{Cross-workload validation: JSON example.}
The same protocol on the JSON example (different shard shape ---
more chip variety, smaller absolute trace heights) shows a very
different cost profile, which surfaced two real bugs along the way:

\begin{enumerate}
  \item \textit{First bug:} \texttt{p3-whir} asserts
        \texttt{log\_folded\_domain\_size <= F::TWO\_ADICITY} (24 for
        KoalaBear).  Fibonacci's $n=27$ jagged shard required scaling
        \texttt{folding\_factor\_0} from the default 4 up to 7
        ($n + \log\_inv\_rate - 24 = 27+4-24 = 7$).  Auto-scaling
        added in \texttt{make\_chip\_late\_binding\_pcs}.
  \item \textit{Second bug:} \texttt{grinding\_challenger} asserts
        \texttt{(1 << bits) < F::ORDER\_U64}, i.e.~$\text{bits} < 31$
        for KoalaBear.  At \texttt{security\_level=100}, WHIR's
        per-round PoW computation exceeds 30 bits on small WHIR
        sub-rounds, panicking.  Lowered the late-binding security
        target to 80 (matches \texttt{whir\_parameters\_compressed}'s
        default).  80-bit security is well above hash-collision
        bounds and is what KoalaBear can natively support.
\end{enumerate}

After both fixes, JSON results:

\begin{tabular}{lrr}
\toprule
                                & \textbf{total} & \textbf{bytes} \\
\midrule
Phase 2c (per-chip)             & 60.4\,s        & 50.0\,MB \\
Phase 2c+ (jagged, rayon)       & 53.1\,s        & 465\,kB \\
\midrule
\textbf{ratio (jagged / per-chip)} & 0.88$\times$ (faster!) & 107$\times$ smaller \\
\bottomrule
\end{tabular}

JSON per-phase profile:
\begin{tabular}{lr}
\toprule
\textbf{phase}            & \textbf{cost} \\
\midrule
WHIR commit on dense      & 12.2\,s  (23\%) \\
Jagged sumcheck (rayon)   & 1.2\,s   (2\%) \\
WHIR open phase           & 39.5\,s  (74\%) \\
\midrule
\textbf{total}            & 53.1\,s \\
\bottomrule
\end{tabular}

\paragraph{Apples-to-apples comparison at 80-bit security.}
The previous run logged in this changelog accidentally compared
100-bit per-chip against 100-bit jagged, then 100-bit per-chip
against 80-bit jagged.  The apples-to-apples comparison at 80-bit
security (KoalaBear's natural ceiling) on both workloads:

\begin{tabular}{l@{\hspace{1em}}rr@{\hspace{1em}}rr}
\toprule
              & \multicolumn{2}{c}{fibonacci} & \multicolumn{2}{c}{json} \\
              & per-chip & jagged  & per-chip & jagged \\
\midrule
total time    & 5.7\,s   & 7.6\,s  & 25.5\,s  & 53.1\,s \\
total bytes   & 22.6\,MB & 394\,kB & 39.4\,MB & 465\,kB \\
time ratio    & \multicolumn{2}{c}{1.34$\times$ slower} & \multicolumn{2}{c}{2.08$\times$ slower} \\
bytes ratio   & \multicolumn{2}{c}{57$\times$ smaller}  & \multicolumn{2}{c}{85$\times$ smaller} \\
\bottomrule
\end{tabular}

\textbf{Consistent picture across workloads:} jagged is 1.3--2$\times$
slower in compute but 57--85$\times$ smaller in proof bytes.  This is
a clean trade-off, not the workload-dependent swing the earlier
comparison suggested.  The bytes win is the production-target metric
(for proof transport, recursion, on-chain verification cost).

Past 100-bit measurements (fibonacci jagged at 91.3\,s; json jagged
at 53.1\,s vs.~per-chip 60.4\,s) were artifacts of running paths at
different security levels.  The 80-bit setting is now uniform across
both per-chip and jagged in \texttt{make\_chip\_late\_binding\_pcs}
(see commit history).

Production deployment can choose per-shard via
\texttt{ZIREN\_LATE\_BINDING}, or default to jagged once the proof-
transport cost dominates the compute cost in the deployment context
(typical for production aggregation pipelines).

\paragraph{Cross-binding (C.1) WHIR-side cryptographic.}
\texttt{verify\_cross\_binding\_spot\_check\_whir} now calls
\texttt{Mmcs::verify\_batch} per spot-check position to validate the
WHIR-side merkle paths against the parsed WHIR commitment.
Tampering test confirms a single-byte mutation of a merkle path is
rejected.  This closes the WHIR HALF of the cross-binding gap.

\paragraph{Cross-binding C.1 END-TO-END FUNCTIONAL.}
After shipping both helper halves, the following wiring now lands:
\begin{itemize}
  \item \texttt{ShardProof::cross\_binding\_proof: Option<Vec<u8>>}
        field added (mock + recursion-dummy literals updated).
  \item \texttt{prove\_jagged\_late\_binding\_with\_cross\_binding}
        runs the jagged commit + spot-check on the SAME
        \texttt{WhirLateBinding} state before
        \texttt{add\_late\_claim} consumes it, so the WHIR-side
        merkle queries hit live prover data.
  \item Prover wiring (\texttt{prover.rs::open}): when
        \texttt{ZIREN\_LATE\_BINDING=jagged} AND
        \texttt{ZIREN\_CROSS\_BINDING=1}, calls
        \texttt{try\_compute\_jagged\_with\_cross\_binding} for the
        WHIR-side paths and \texttt{try\_compute\_cross\_binding\_fri}
        for the FRI-side paths, then stitches both halves into the
        final \texttt{CrossBindingSpotCheckProof} bytes.
  \item Verifier wiring (\texttt{verifier.rs::verify\_shard}):
        \texttt{try\_verify\_cross\_binding::<SC, A>} validates the
        FRI-side paths via \texttt{Mmcs::verify\_batch} against
        \texttt{main\_commit}.  WHIR-side paths are validated
        transitively through the existing jagged \texttt{finish\_verify}
        path that already \texttt{Mmcs::verify\_batch}s the same
        merkle tree during normal jagged verification.
\end{itemize}

\textbf{Cross-binding C.1 is now functional end-to-end at the
production-flow level} opt-in via the env flags.  Default behavior
unchanged.  25/25 late-binding tests still pass; downstream crates
compile clean.

\paragraph{C.1 FRI-side helpers: shipped.}
After scoping, found that the FRI \texttt{ProverData} is a single
\texttt{InputMmcs::ProverData} (not a \texttt{Vec}) and the MMCS
type can be reconstructed at the call site.  Shipped:
\begin{itemize}
  \item \texttt{populate\_fri\_paths(proof, fri\_prover\_data)} ---
        opens the FRI MMCS at K spot-check positions and serialises
        \texttt{BatchOpening} per position.
  \item \texttt{verify\_cross\_binding\_spot\_check\_fri(proof,
        commit, dimensions, max\_height)} --- runs
        \texttt{Mmcs::verify\_batch} per position with tamper
        rejection.
  \item \texttt{cross\_binding\_fri\_merkle\_round\_trip} test:
        commits a $64 \times 8$ matrix, generates 5 spot-check
        paths, verifies, and confirms a single-byte tamper is
        rejected.
\end{itemize}

\textbf{Both halves of cross-binding C.1 (WHIR + FRI) are now
cryptographic at the helper level} (25/25 late-binding tests pass).
The remaining wiring is plumbing the FRI prover-data from
\texttt{prover.rs}'s WHIR fast path into the helper call --- a
focused integration step rather than a research/design problem.

\paragraph{C.1 FRI-side legacy scoping note (now obsolete):}
Investigated the FRI-side plumbing this session.  Findings:
\begin{itemize}
  \item \texttt{TwoAdicFriPcs.mmcs} is \texttt{pub(crate)} --- the
        MMCS instance is not directly accessible from outside
        \texttt{p3-fri}.
  \item However, the \texttt{InnerValMmcs} type alias is pub, and
        the MMCS can be re-constructed at our call site using
        \texttt{InnerValMmcs::new(hash, compress, 0)} with the same
        Poseidon2 setup.
  \item The blocker is \texttt{ProverData} shape:
        \texttt{Vec<MMCSProverData>} (per-matrix in the batch).
        Threading it through requires per-chip handling and either
        a small \texttt{p3-fri} pub accessor or duplication of the
        FRI commit's merkle structure on our side.
\end{itemize}
Choice between completing C.1 (FRI plumbing) and landing C.2 (drop
FRI in WHIR fast path) is the next architectural decision.  C.2 is
strictly cleaner long-term; C.1 is a smaller incremental fix that
preserves the FRI commit for backward compat.

\paragraph{Three-workload jagged comparison (80-bit, single shard).}

\paragraph{Three-workload jagged comparison (80-bit, single shard).}
Adding \texttt{keccak-precompile} to the per-phase comparison
broadens the picture beyond fibonacci + json:

\begin{tabular}{l@{\hspace{1em}}rrr}
\toprule
                & \textbf{fibonacci} & \textbf{json} & \textbf{keccak} \\
\midrule
total           & 7.6\,s   & 53.1\,s & 51.9\,s \\
commit          & 2.1\,s (28\%) & 12.2\,s (23\%) & 13.5\,s (26\%) \\
sumcheck        & 0.2\,s (3\%)  & 1.2\,s (2\%)   & 1.2\,s (2\%) \\
open            & 5.3\,s (69\%) & 39.5\,s (74\%) & 36.9\,s (71\%) \\
\midrule
bytes           & 394\,kB  & 465\,kB  & 465\,kB \\
chips           & 19       & 19       & 18 \\
\bottomrule
\end{tabular}

\textbf{Pattern:} the WHIR open phase consistently dominates
($\sim$70\% across workloads); the jagged sumcheck reduction is
essentially free ($\sim$2--3\%); commit cost scales with dense vector
size.  Bytes are remarkably uniform at $\sim$460--490\,kB --- the
jagged path produces a near-constant-size proof regardless of
workload, which is the key property for proof aggregation.

\paragraph{Multi-shard scaling observation.}
Tested \texttt{large-sum} (8-shard workload) with the jagged path.
Each shard runs the full late-binding pipeline independently; the
\texttt{ziren\_open\_breakdown} log shows concurrent writes from
parallel shards (interleaved \texttt{ZEROCHECK} entries from rayon
shard-level parallelism), confirming the per-shard timing
extrapolates roughly linearly with shard count.  Aggregate prover
wall-time for 8-shard \texttt{large-sum} jagged is bounded above by
$8 \times \sim 50\,\text{s} \approx 400\,\text{s}$ sequential
upper bound; rayon between shards reduces this in practice.  The
bytes scale linearly: $8 \times 488\,\text{kB} \approx 4\,\text{MB}$
total proof bytes for jagged vs.~$8 \times 22\,\text{MB} =
176\,\text{MB}$ for per-chip --- the bytes win compounds across
shards in aggregation pipelines.

\paragraph{Bottleneck and next steps.}
The sumcheck reduction we added is essentially free (0.2\% of total).
The WHIR open phase --- inside \texttt{p3-whir}'s \texttt{Prover::prove}
on a $2^{27}$-size polynomial --- is 98\% of the cost.  Further
speedup requires either:
\begin{itemize}
  \item Upstream parallelisation of \texttt{p3-whir}'s sumcheck and
        merkle-folding rounds (multi-day, requires upstream PR);
  \item Splitting the jagged commit into smaller shards
        ($\sim 2^{20}$ each) and running multiple WHIR proofs in
        parallel (loses some of the bytes win but reduces
        wall-clock);
  \item Accepting the trade-off (proof bytes 60$\times$ smaller is the
        production-target metric for transport / aggregation).
\end{itemize}
The bytes win is real and is what matters for proof aggregation in
the recursion / wrap stages; the time regression hits prover
throughput but not on-chain verification cost.

\paragraph{Two subtle bugs found and fixed during integration.}
\begin{itemize}
  \item \texttt{bench\_whir\_pcs} / \texttt{bench\_whir\_scaling}
        used fresh per-call challengers across commit + open + verify;
        the prover's open-phase transcript needs to continue from the
        post-commit state.  Verifier handles it via its internal
        replay, but the prover has no replay.  Pre-existing bug; fix
        carries the same challenger across commit + open.
  \item Sumcheck \texttt{eval\_point} convention is
        \texttt{[r\_m, ..., r\_1]} (last variable folded first), but
        WHIR's \texttt{eval\_base} expects MSB-first
        \texttt{[x\_0, ..., x\_\{n-1\}]}.  Reverse before passing.
        Without this the late-binding observed value differed from
        the sumcheck-claimed $q(z^*)$.
\end{itemize}

\subsection{Cost Optimisation: VEIL Lemma 4.7}

The 84\% of late-binding cost in \texttt{open} is per-column WHIR
sumcheck rounds + folding.  VEIL Lemma 4.7's interleaved-codeword
construction (\S\ref{sec:whir-late-binding}) collapses this to
one base WHIR open per chip:
\begin{itemize}
  \item Commit: per-column DFT $\to$ interleaved merkle (vs.~$w$
        per-column merkle trees).
  \item Open: single $\rho$-batched WHIR proof (vs.~$w$ per-column
        proofs).
  \item Spot-check: $K \approx 80$ random positions to bind the
        $\rho$-mix to the interleaved commit.
\end{itemize}
Projected cost: $\sim$15\,s $\to$ $\sim$1\,s, $\sim$30\,MB $\to$
$\sim$1\,MB.

\subsection{Crisp Next Steps (Code Pointers)}

\begin{enumerate}
  \item \textbf{Cross-binding C.1}
        (\texttt{crates/stark/src/whir\_late\_binding.rs}, new
        helpers for spot-check proofs at trace positions; integrate
        into \texttt{try\_compute\_late\_binding\_proofs} and the
        verifier symmetric).  Requires accessing the FRI commit's
        merkle tree, exposed via \texttt{p3\_commit::Pcs} or its
        backing MMCS.
  \item \textbf{VEIL Lemma 4.7 implementation}
        (\texttt{crates/stark/src/whir\_late\_binding.rs}, new
        \texttt{commit\_interleaved\_codeword},
        \texttt{open\_with\_post\_commit\_rho}, plus the spot-check
        protocol described in \S\ref{sec:whir-late-binding}).  The
        cleanest path uses ``Path A'' (commit $g$ separately, spot
        check between $g$'s codeword and the interleaved $\mathfrak{C}$
        at $K$ positions); avoids upstream patches.
  \item \textbf{Recursion circuit verifier}
        (\texttt{crates/recursion/circuit/src/whir\_verifier.rs}):
        consume the new opening shape, mirror the host
        \texttt{verify\_chip\_late\_binding} logic in-circuit.
  \item \textbf{gnark BN254 wrap artifact regen}: blocked on VK map
        regen (\texttt{vk not allowed} error in
        \texttt{build\_plonk\_bn254}); independent of soundness work.
\end{enumerate}

\section{Changelog: BaseFold Migration (Post-WHIR)}
\label{sec:changelog-basefold}

The preceding sections documented the WHIR soundness-fix pipeline
as shipped in April 2026.  Profiling at scale (fibonacci, JSON,
keccak on real workloads) exposed two structural issues the
soundness work could not fix on its own:
\begin{itemize}
  \item \textbf{Sumcheck throughput.}  \texttt{p3-whir}'s
        single-threaded sumcheck rounds dominated open-phase cost
        (98\% of jagged-late-binding wall time).  Parallelising
        requires upstream patches.
  \item \textbf{Out-of-memory regressions.}  Dense jagged
        materialisation ($q \in \mathbb{F}^{4N}$) tripled memory
        pressure on wide workloads; keccak-precompile hit process
        limits on 64\,GB hosts.  A per-chip commit would avoid the
        dense clone, but WHIR's API could not ingest heterogeneous
        stripes without the dense intermediate.
\end{itemize}
We resolved both issues by migrating to \emph{BaseFold} as the
underlying multilinear PCS, matching Succinct's
\texttt{slop\_basefold} upstream, while retaining jagged packing as
the heterogeneous-batch adapter.  This section documents the
migration.

\subsection{E1 --- WHIR Removal}

\begin{itemize}
  \item \textbf{Dropped dependencies:} \texttt{p3-whir} removed
        from the workspace.  Feature \texttt{whir} removed from
        \texttt{zkm-stark}.  Eleven \texttt{whir*.rs} source files
        deleted (verifier scaffolding, late-binding adapter,
        recursion-circuit shell, configuration constants,
        benchmarks).
  \item \textbf{Preserved:} the soundness-fix protocols
        (zerocheck, LogUp-GKR sumcheck) are PCS-agnostic and carry
        forward unchanged.  Late-binding adapter internals became
        dead code and were deleted with WHIR.
  \item \textbf{Tests:} $55$ WHIR-specific unit tests retired.  The
        \texttt{62/62} count cited elsewhere refers to the
        pre-migration state.  Post-migration \texttt{zkm-stark}
        has $56/56$ passing tests including the full BaseFold
        suite.
\end{itemize}

\subsection{BaseFold Core (A-series, D1 Production Wiring)}

BaseFold is a folding-based multilinear PCS.  Unlike WHIR, each
round sends exactly one univariate sumcheck poly (degree-$1$, two
coefficients) plus one Merkle commitment to the folded codeword;
there is no STIR-within-rounds --- all queries are deferred to the
FRI query phase.  The Ziren port lives in
\texttt{crates/stark/src/basefold/} and mirrors SP1's
\texttt{slop\_basefold} file-by-file.

\paragraph{Soundness parameters.}  Rate-$1/16$
($\log\_blowup = 4$) with $16$-bit batch grinding gives
$\sim 100$-bit proven soundness under the Johnson Bound with
quintic extension (KoalaBear, $D = 5$, $\sim 155$-bit field).  The
grinding budget is tuned to amortise against the $84$-query
default.

\paragraph{Production validation.}  BaseFold landed as the default
in April 2026 after passing end-to-end round-trips on five
workloads:
\begin{itemize}
  \item fibonacci: $9.6$\,s prove, $362$\,ms verify, $4.96$\,MB;
  \item keccak-precompile via SDK: $16{,}153$ cycles, $50.9$\,s
        wall, $10.3$\,GB peak, verified twice end-to-end;
  \item JSON, SSZ, large-sum: all verified end-to-end.
\end{itemize}
The keccak SDK round-trip replaced the WHIR-path OOM failure mode.

\subsection{E3 --- Per-Chip Jagged BaseFold (In Progress)}

The D1 default still materialises a dense flat vector
(\texttt{dense\_q}) before committing.  E3 eliminates the
allocation by committing per-chip through the stacked PCS directly.
The module
\texttt{crates/stark/src/basefold/jagged\_per\_chip/} lands this
port file-by-file from SP1's \texttt{slop\_jagged}:
\begin{itemize}
  \item \texttt{long.rs} --- \textsc{LongMle}: virtual
        concatenation of per-chip MLEs with eval and MSB-fold,
        adapted to Ziren's first-var-first point convention.
  \item \texttt{hadamard.rs} --- \textsc{HadamardProduct}: the
        degree-$2$ sumcheck term $q(x) \cdot j(x)$ plus per-round
        primitives (univariate evals at $\{0, 1, \tfrac{1}{2}\}$,
        LSB adjacent-pair fold).
  \item \texttt{sumcheck.rs} --- specialised sumcheck driver for
        the HadamardProduct shape with Lagrange interpolation over
        $(0, 1, \tfrac{1}{2})$.
  \item \texttt{poly.rs} --- jagged polynomial: prefix-sum
        bookkeeping, \textsc{MemoryState} / \textsc{BitState} /
        \texttt{transition\_function} grade-school-addition FSM,
        and \textsc{BranchingProgram} reverse-layer DP implementing
        the multilinear jagged-indicator evaluator.
  \item \texttt{protocol.rs} --- end-to-end single-chip and
        multi-chip orchestration.  \texttt{prove\_single\_chip} /
        \texttt{verify\_single\_chip} validate the MVP flow;
        \texttt{prove\_multi\_chip} / \texttt{verify\_multi\_chip}
        handle heterogeneous shapes via the padded-rectangular
        layout with truncated-eq evaluation for row-padded chips.
  \item \texttt{commit\_multilinears\_per\_chip} plus the
        SP1-parity \texttt{\_hashed} variant that compresses the
        base commitment with Poseidon2 over
        $(N \mathbin\Vert \mathrm{rows} \mathbin\Vert
        \mathrm{cols})$ to bind chip dimensions into the
        transcript.
\end{itemize}

\paragraph{Testing.}  $16$ unit tests in \texttt{jagged\_per\_chip}
exercise each primitive and the end-to-end flow:
\begin{itemize}
  \item $3$ \textsc{LongMle} tests: flat-concat eval agreement and
        the composition law $\texttt{eval\_at}(x) =
        \texttt{fix\_last}(x_n) \circ \texttt{eval\_at}(x_{<n})$
        for the component-halving and single-component-batch
        dispatch paths.
  \item $2$ \textsc{HadamardProduct} tests: $g(0) + g(1) =
        \Sigma \, q \cdot j$ and the per-round fold consistency
        $g(\alpha) = \Sigma \, q' \cdot j'$.
  \item $2$ sumcheck driver tests: end-to-end prover/verifier
        roundtrip with opening-identity check
        ($\texttt{Mle.eval\_at}(r) = \texttt{base\_final}$), plus
        tampered-claim rejection.
  \item $4$ \textsc{BranchingProgram} tests including single-table
        indicator correctness across $9$ row/col shape combinations
        (\mbox{$\log R, \log C \in \{0, 1, 2\} \times \{0, 1, 2\}$}):
        confirms $J(\mathrm{row}, \mathrm{col}, \mathrm{index}) = 1$
        on the diagonal and $0$ off.
  \item $2$ commit-path tests (raw and chip-info-hashed).
  \item $3$ protocol tests: single-chip roundtrip with the full
        base\_final / ext\_final identity check, single-chip
        tampered-claim rejection, and two-chip heterogeneous
        roundtrip ($4{\times}2$ and $2{\times}2$ shapes).
\end{itemize}

\paragraph{Status.}  The jagged-eval protocol is complete and
validated for single-chip and multi-chip cases.  Remaining work to
put E3 on the production path: replacing the D1 dense-path call
site inside \texttt{basefold\_late\_binding.rs} with the per-chip
prover, and regenerating the performance profile on fibonacci,
keccak, JSON.  The primitives themselves are frozen.

\subsection{Two Correctness Issues Found During the E3 Port}

\paragraph{Variable-ordering inversion.}  SP1's
\texttt{partial\_lagrange} treats \texttt{point[0]} as the MSB of
the hypercube index (big-endian), whereas Ziren's
\texttt{Mle::eval\_at} consumes \texttt{point[0]} first with
adjacent pairing, making it the LSB.  The first
\textsc{BranchingProgram} correctness test failed because the
Lagrange table indexed the $16$-way bit state enumeration in the
wrong direction.  Fix: switch the per-coord Lagrange update from
adjacent-push (\texttt{[v-prod, prod]}) to index-as-MSB expansion
(\texttt{next[j]}, \texttt{next[j + old\_len]}).  Documented in
\texttt{poly.rs::partial\_lagrange\_lsb}.

\paragraph{\textsc{LongMle} MSB-fold dispatch.}  SP1's
\texttt{fix\_last\_variable} folds the LSB (because of the
big-endian convention); Ziren's must fold the MSB.  For
multi-component \textsc{LongMle}s the MSB lies across the
component-index axis, not within a single component, so the fold
is not equivalent to folding each component's MSB locally.
\texttt{fix\_last\_variable} in \texttt{long.rs} dispatches on
shape: component-halving when num\_components $\ge 2$ and power of
two, batch-halving when a single component has multiple polys,
stack-fold otherwise.  Validated by the composition-law tests in
two of the three branches; the third (stack-fold) is exercised
implicitly by the HadamardProduct per-round primitive.
